\begin{afterword}
\summaryhead

The post opens with a short history. Hugh Everett proposed many-worlds in 1957; the paper was ignored, Bohr did not take him seriously, and Everett left physics for defense work and business. Physicists at large heard of the idea only through Bryce DeWitt's article of 1970. The post then asks whether the order of discovery decides which interpretation is popular, and imagines a world where many-worlds came first and no one proposed a collapse until 1957.

In that world Huve Erett brings Biels Nohr a paper saying that at a measurement all of the wavefunction but one part vanishes. Nohr lists the ways such a collapse would be unique in physics: non-linear, non-local, at odds with relativity, informally specified, untestable as stated, and supported by no evidence. Each time, Huve patches the theory to fit. At the end Huve imagines a world where collapse came first and no one thought of many worlds for thirty years, and Nohr laughs at the idea. That world is ours.

\responsehead

The question is a fair one: would a theory be accepted if it were proposed second rather than first? The history the post gives is right in outline: the 1957 paper, its neglect, Wheeler sending Everett to Bohr, DeWitt's article in 1970. One link is out of order. Everett had already taken a job at the Pentagon in 1956, to keep his draft deferment, before the paper appeared and three years before he met Bohr; he did not leave physics because Bohr rejected him.

The satire's force depends on the collapse theory it picks. Huve's is the textbook postulate, and Nohr's strongest charges fit it: it is informal, it cannot say when or in which basis the collapse happens, and so it cannot be tested. Those charges do not fit the collapse theories physicists had published by 2008. The GRW theory of 1986 fixes the basis, a rate and a width, and such theories are being tested; the simplest Diósi--Penrose model has been ruled out. None of these tests has found a collapse, which supports the narrower charge Nohr also makes, that there is no evidence for one. But the story sets many-worlds against the weakest form of its rival, and does not mention the stronger ones.

The hard questions also go to one side. Nohr admits that his side does not know why the Born statistics hold. The matching problem for many-worlds, how there can be probabilities when every outcome occurs, is called its most serious difficulty in the \textsc{Stanford Encyclopedia}, and Huve does not raise it. A thought experiment of this kind can show that familiarity may shape which theory looks absurd. It cannot show which arguments are sound, since the arguments are the same whichever theory came first. The fuller case for many-worlds is made later in the sequence.

The ending mocks our world's founders for never asking whether the laws were universal. They asked. Von Neumann in 1932 let the measuring instrument follow the Schrödinger equation and kept the collapse at the observer. Schrödinger in 1935 carried the superposition up to a cat and called the result ``quite ridiculous.'' They judged the answer absurd because observation shows one outcome; that is a reason, whatever its worth, not a failure to think of the question. The last line replaces argument with a jab about cryonics and a president. Smaller slips: the Wheeler--DeWitt equation, cited as the well-tested law, is a proposed equation of quantum gravity, and Everett did not invent the use of Lagrange multipliers in optimization.

In short: a fair question about whether the order of discovery shapes belief, answered with a satire that sets many-worlds against the textbook collapse postulate rather than the specified, testable collapse theories, and that mocks the founders for not asking a question they did ask.
\end{afterword}
